\documentclass[10pt,a4paper]{amsart}
\usepackage{amsmath,amssymb,amsthm,graphicx}
\usepackage{verbatim}
\newenvironment{reference}{\comment}{\endcomment}
%\newenvironment{reference}{}{}
\newenvironment{slogan}{\comment}{\endcomment}
\newenvironment{history}{\comment}{\endcomment}

% For commutative diagrams we use Xy-pic
\usepackage[all]{xy}

% We use 2cell for 2-commutative diagrams.
\xyoption{2cell}
% Theorem environments.
%
\theoremstyle{plain}
\newtheorem{theorem}[subsection]{Theorem}
\newtheorem{proposition}[subsection]{Proposition}
\newtheorem{lemma}[subsection]{Lemma}

\theoremstyle{definition}
\newtheorem{definition}[subsection]{Definition}
\newtheorem{example}[subsection]{Example}
\newtheorem{exercise}[subsection]{Exercise}
\newtheorem{situation}[subsection]{Situation}

\theoremstyle{remark}
\newtheorem{remark}[subsection]{Remark}
\newtheorem{remarks}[subsection]{Remarks}

\numberwithin{equation}{subsection}

% Macros
%
\def\lim{\mathop{\mathrm{lim}}\nolimits}
\def\colim{\mathop{\mathrm{colim}}\nolimits}
\def\Spec{\mathop{\mathrm{Spec}}}
\def\Hom{\mathop{\mathrm{Hom}}\nolimits}
\def\Ext{\mathop{\mathrm{Ext}}\nolimits}
\def\SheafHom{\mathop{\mathcal{H}\!\mathit{om}}\nolimits}
\def\SheafExt{\mathop{\mathcal{E}\!\mathit{xt}}\nolimits}
\def\Sch{\mathit{Sch}}
\def\Mor{\mathop{\mathrm{Mor}}\nolimits}
\def\Ob{\mathop{\mathrm{Ob}}\nolimits}
\def\Sh{\mathop{\mathit{Sh}}\nolimits}
\def\NL{\mathop{N\!L}\nolimits}
\def\CH{\mathop{\mathrm{CH}}\nolimits}
\def\proetale{{pro\text{-}\acute{e}tale}}
\def\etale{{\acute{e}tale}}
\def\QCoh{\mathit{QCoh}}
\def\Ker{\mathop{\mathrm{Ker}}}
\def\Im{\mathop{\mathrm{Im}}}
\def\Coker{\mathop{\mathrm{Coker}}}
\def\Coim{\mathop{\mathrm{Coim}}}

% Boxtimes
%
\DeclareMathSymbol{\boxtimes}{\mathbin}{AMSa}{"02}

%
% Macros for moduli stacks/spaces
%
\def\QCohstack{\mathcal{QC}\!\mathit{oh}}
\def\Cohstack{\mathcal{C}\!\mathit{oh}}
\def\Spacesstack{\mathcal{S}\!\mathit{paces}}
\def\Quotfunctor{\mathrm{Quot}}
\def\Hilbfunctor{\mathrm{Hilb}}
\def\Curvesstack{\mathcal{C}\!\mathit{urves}}
\def\Polarizedstack{\mathcal{P}\!\mathit{olarized}}
\def\Complexesstack{\mathcal{C}\!\mathit{omplexes}}
% \Pic is the operator that assigns to X its picard group, usage \Pic(X)
% \Picardstack_{X/B} denotes the Picard stack of X over B
% \Picardfunctor_{X/B} denotes the Picard functor of X over B
\def\Pic{\mathop{\mathrm{Pic}}\nolimits}
\def\Picardstack{\mathcal{P}\!\mathit{ic}}
\def\Picardfunctor{\mathrm{Pic}}
\def\Deformationcategory{\mathcal{D}\!\mathit{ef}}
\usepackage{cleveref}
\crefname{originaltheorem}{original theorem}{original theorems}
\Crefname{originaltheorem}{Original theorem}{Original theorems}
\makeatletter
\newlabel{categories-theorem-adjoint-functor}{{0AHQ}{source}{}{sourcepointer-1}{}}
\newlabel{categories-theorem-adjoint-functor@cref}{{[originaltheorem][0][]0AHQ}{[1][0][]source}{}{}{}}
\makeatother
\crefname{originallemma}{original lemma}{original lemmas}
\Crefname{originallemma}{Original lemma}{Original lemmas}
\makeatletter
\newlabel{lemma-set-iso-classes-bounded-size}{{079E}{source}{}{sourcepointer-2}{}}
\newlabel{lemma-set-iso-classes-bounded-size@cref}{{[originallemma][0][]079E}{[1][0][]source}{}{}{}}
\makeatother
\title{The Gabriel-Popescu functor has a left adjoint}
\author{The Stacks Project Authors}
\date{stacks-project:injectives.tex:local-snapshot-20261012; GNU FDL 1.2 or later}
\begin{document}
\maketitle
\noindent\textit{Editorial note.} All delivered mathematical text is the author's own, copied from the pinned Stacks Project chapter file sites-cohomology.tex; no statement or proof step was rewritten, reordered or removed.\par\medskip
\noindent\textit{Editorial note.} The proof is the author's complete argument, containing two independent routes to the conclusion. The single omitted step is the sentence 'and is functorial; we omit the details.' concerning the routine functoriality of the explicitly constructed assignment, which follows from the representatives chosen in the construction; it is preserved verbatim and the rest of both arguments is written out in the counted span.\par\medskip
\noindent\textit{Editorial note.} Every labelled reference occurring inside the counted spans is carried in assembly.external\_references as a non-load-bearing navigation pointer (same-chapter ones at their definition site, cross-chapter ones at the invocation inside the counted span), so no load-bearing mathematics is replaced by a stub. No bibliography citation occurs inside the counted spans.\par\medskip
\noindent\textit{Editorial note.} The same-name check for this card scanned every staged Stacks chapter of this repository -- all sources/stacks-*/pinned-original-source.tex files, the card's own chapter excluded -- and the labels found are recorded in dedup.same\_name\_cross\_chapter on the card.\par\medskip
\noindent\textit{Editorial note.} The preamble is the chapter's own retained minimal preamble (sources/stacks-sites-cohomology/retained-preamble.tex), a verbatim excerpt of the native preamble: the Xy-pic setup and the native theorem declarations with the mathematical macros.\par\medskip
\noindent\textit{Allowed standard prerequisites (editorial).}\begin{itemize}
\item The Gabriel-Popescu functor. Same chapter, the construction of the functor G from abelian groups to the Grothendieck abelian category and its properties established before this lemma. Supplies the functor whose left adjoint is constructed and the properties of it used in both proofs.
\item The adjoint functor theorem. Categories, the adjoint functor theorem for locally presentable or well-powered categories. Supplies the first proof, in which the existence of the left adjoint follows from the preservation of limits and the solution set condition.
\item The explicit construction of the left adjoint. Same chapter, the construction of the functor F on modules by generators and relations. Supplies the second proof, where the author builds the left adjoint from copies of the functor and checks the adjunction on Hom sets.
\end{itemize}\medskip
\ifdefined\hypertarget\hypertarget{sourcepointer-1}{}\fi
\noindent\textit{Source reference.} Original theorem 0AHQ: sources/stacks-categories/pinned-original-source.tex, source lines 3964--3964. Navigation pointer invoked at proof line 2773; the counted span contains the pointer only and the cited result is not reproduced inside it, so the pointer is registered as a non-load-bearing navigation stub in the established form for Stacks material. No bibliography citation is involved.\par\medskip
\ifdefined\hypertarget\hypertarget{sourcepointer-2}{}\fi
\noindent\textit{Source reference.} Original lemma 079E: sources/stacks-injectives/pinned-original-source.tex, source lines 1616--1616. Navigation pointer invoked at proof line 2780; the counted span contains the pointer only and the cited result is not reproduced inside it, so the pointer is registered as a non-load-bearing navigation stub in the established form for Stacks material. No bibliography citation is involved.\par\medskip
\begin{lemma}
\label{lemma-gabriel-popescu-left-adjoint}
The functor $G$ above has a left adjoint
$F : \text{Mod}_R \to \mathcal{A}$.
\end{lemma}

\begin{proof}
We will give two proofs of this lemma.

\medskip\noindent
The first proof will use the adjoint functor theorem, see
Categories, Theorem \ref{categories-theorem-adjoint-functor}.
Observe that $G : \mathcal{A} \to \text{Mod}_R$ is left exact and sends
products to products. Hence $G$ commutes with limits. To check the set
theoretical condition in the theorem, suppose that $M$ is an object of
$\text{Mod}_R$. Choose a suitably large cardinal $\kappa$ and denote $E$
a set of objects of $\mathcal{A}$ such that every object $A$ with
$|A| \leq \kappa$ is isomorphic to an element of $E$. This is possible
by Lemma \ref{lemma-set-iso-classes-bounded-size}. Set
$I = \coprod_{A \in E} \Hom_R(M, G(A))$.
We think of an element $i \in I$ as a pair $(A_i, f_i)$.
Finally, let $A$ be an arbitrary object of $\mathcal{A}$
and $f : M \to G(A)$ arbitrary. We are going to think of
elements of $\Im(f) \subset G(A) = \Hom_\mathcal{A}(U, A)$
as maps $u : U \to A$. Set
$$
A' = \Im(\bigoplus\nolimits_{u \in \Im(f)} U \xrightarrow{u} A)
$$
Since $G$ is left exact, we see that $G(A') \subset G(A)$
contains $\Im(f)$ and we get $f' : M \to G(A')$ factoring $f$.
On the other hand, the object $A'$ is
the quotient of a direct sum of at most $|M|$ copies of $U$.
Hence if $\kappa = |\bigoplus_{|M|} U|$, then we see that $(A', f')$
is isomorphic to an element $(A_i, f_i)$ of $E$ and we conclude that $f$
factors as $M \xrightarrow{f_i} G(A_i) \to G(A)$ as desired.

\medskip\noindent
The second proof will give a construction of $F$ which will show
that ``$F(M) = M \otimes_R U$'' in some sense. Namely, for any
$R$-module $M$ we can choose a resolution
$$
\bigoplus\nolimits_{j \in J} R \to
\bigoplus\nolimits_{i \in I} R \to
M \to 0
$$
Then we define $F(M)$ by the corresponding exact sequence
$$
\bigoplus\nolimits_{j \in J} U \to
\bigoplus\nolimits_{i \in I} U \to
F(M) \to 0
$$
This construction is independent of the choice of the resolution
and is functorial; we omit the details.
For any $A$ in $\mathcal{A}$ we obtain an exact sequence
$$
0 \to \Hom_\mathcal{A}(F(M), A) \to
\prod\nolimits_{i \in I} G(A) \to
\prod\nolimits_{j \in J} G(A)
$$
which is isomorphic to the sequence
$$
0 \to \Hom_R(M, G(A)) \to
\Hom_R(\bigoplus\nolimits_{i \in I} R, G(A)) \to
\Hom_R(\bigoplus\nolimits_{j \in J} R, G(A))
$$
which shows that $F$ is the left adjoint to $G$.
\end{proof}

\end{document}
